\chapter{Fonctions Linéaires, Affines \& Graphiques}
\label{chap:fonctions_affines}

% ==============================================================================
% 1. OBJECTIFS & COMPÉTENCES DU RÉFÉRENTIEL
% ==============================================================================
\begin{cadreretenu}[Objectifs \& Compétences du DNB Pro]
\begin{itemize}
    \item \textbf{Connaissances} : Notion de fonction, variable, image, antécédent, fonction linéaire $f(x)=ax$, fonction affine $f(x)=ax+b$, coefficient directeur, ordonnée à l'origine.
    \item \textbf{Capacités} :
    \begin{itemize}
        \item \capRealiser\ Calculer l'image d'un nombre par une fonction définie par une formule.
        \item \capRealiser\ Déterminer un antécédent en résolvant une équation du premier degré.
        \item \capRealiser\ Représenter graphiquement une fonction linéaire (droite passant par l'origine) ou affine.
        \item \capAnalyser\ Lire graphiquement des images et des antécédents, déterminer le coefficient directeur et l'ordonnée à l'origine.
        \item \capValider\ Modéliser une situation de proportionnalité ou d'abonnement / coût fixe par une fonction.
    \end{itemize}
\end{itemize}
\end{cadreretenu}

\vspace{0.3cm}

% ==============================================================================
% 2. COURS ET NOTIONS ESSENTIELLES
% ==============================================================================
\section{Cours : Notions Fondamentales \& Représentations}

\subsection{Vocabulaire des Fonctions : Image et Antécédent}

\begin{cadredef}[Fonction, Image, Antécédent]
Une \textbf{fonction} $f$ est un processus mathématique qui, à un nombre $x$ de départ (l'antécédent), associe un unique nombre d'arrivée noté $f(x)$ (l'image).
\[ x \overset{f}{\longmapsto} f(x) \]
\begin{itemize}
    \item Un nombre possède \textbf{une seule image}.
    \item Un nombre peut posséder \textbf{aucun, un ou plusieurs antécédents}.
\end{itemize}
\end{cadredef}

\begin{cadreexemple}[Calcul d'image et d'antécédent]
Soit la fonction $f$ définie par $f(x) = 3x - 5$.
\begin{itemize}
    \item \textbf{Calculer l'image de 4} : On remplace $x$ par $4$ : $f(4) = 3 \times 4 - 5 = 12 - 5 = \repbox{7}$.
    \item \textbf{Déterminer l'antécédent de 16} : On résout $f(x) = 16 \iff 3x - 5 = 16 \iff 3x = 21 \iff x = \repbox{7}$.
\end{itemize}
\end{cadreexemple}

\subsection{Fonctions Linéaires \& Fonctions Affines}

\begin{cadreprop}[Caractérisation des fonctions usuelles]
\begin{itemize}
    \item \textbf{Fonction Linéaire} : $f(x) = ax$. Sa courbe représentative est une \textbf{droite qui passe par l'origine} du repère $(0;0)$. Elle modélise une situation de \textbf{proportionnalité}.
    \item \textbf{Fonction Affine} : $f(x) = ax + b$. Sa courbe représentative est une \textbf{droite} qui coupe l'axe des ordonnées au point $(0\,;\,b)$.
    \begin{itemize}
        \item $a$ est le \textbf{coefficient directeur} (ou pente de la droite). Si $a > 0$, la fonction est croissante ; si $a < 0$, elle est décroissante.
        \item $b$ est l'\textbf{ordonnée à l'origine}.
    \end{itemize}
\end{itemize}
\end{cadreprop}

\begin{center}
\begin{tikzpicture}[scale=0.85]
    % Grille et repère
    \draw[step=1cm, slate200, very thin] (-3,-2) grid (5,4);
    \draw[thick, ->, >=stealth, slate700] (-3,0) -- (5,0) node[right] {$x$};
    \draw[thick, ->, >=stealth, slate700] (0,-2) -- (0,4) node[above] {$y$};
    \node[below left, font=\footnotesize] at (0,0) {$O$};

    % Droite affine f(x) = 0.8x + 1
    \draw[line width=1.4pt, colPrimary] (-3,-1.4) -- (3.75,4) node[above right] {$d_1 : f(x) = 0{,}8x + 1$};
    \fill[colPrimary] (0,1) circle (2.5pt) node[left=2pt, font=\footnotesize\bfseries] {$b = 1$};

    % Pente triangle
    \draw[thick, colSecondary, dashed] (0,1) -- (2,1) node[midway, below] {$+2$};
    \draw[thick, colSecondary, dashed] (2,1) -- (2,2.6) node[midway, right] {$+1{,}6$};

    % Droite linéaire g(x) = -x
    \draw[line width=1.4pt, colProp] (-3,3) -- (3,-3) node[below right] {$d_2 : g(x) = -x$};
    \fill[colProp] (0,0) circle (2.5pt);
\end{tikzpicture}
\end{center}

% ==============================================================================
% 3. AUTOMATISMES & QUESTIONS FLASH
% ==============================================================================
\section{Questions Flash / Automatismes}

\begin{automatisme}[Lecture Graphique \& Calculs Rapides (10 min)]
Soit la fonction affine $g(x) = -2x + 6$.
\noindent
\begin{tabularx}{\linewidth}{X|X}
\textbf{1.} Le coefficient directeur est $a = \pointille[1.5cm]$ & \textbf{4.} L'image de $3$ est $g(3) = \pointille[2cm]$ \\
\textbf{2.} L'ordonnée à l'origine est $b = \pointille[1.5cm]$ & \textbf{5.} L'antécédent de $0$ est : $x = \pointille[2cm]$ \\
\textbf{3.} La fonction $g$ est : (Croissante / Décroissante) ? \pointille[2cm] & \textbf{6.} La droite passe par le point $(0\,;\,\pointille[1cm])$. \\
\end{tabularx}
\end{automatisme}

% ==============================================================================
% 4. TRAVAUX DIRIGÉS (TD)
% ==============================================================================
\section{Travaux Dirigés (TD) : Exercices d'Entraînement}

\begin{exercice}[3 pts]{\capRealiser}{Calculs d'images et d'antécédents}
On considère la fonction $f$ définie sur $\mathbb{R}$ par $f(x) = 4x - 7$.
\begin{enumerate}
    \item Calculer $f(2)$, $f(-3)$ et $f(0)$.
    \item Déterminer l'antécédent de $13$ par la fonction $f$.
    \item Le point $M(5\,;\,13)$ appartient-il à la droite représentative de $f$ ? Justifier.
\end{enumerate}
\lignesreponse[5]
\end{exercice}

\begin{exercice}[4 pts]{\capRealiser}{Tracé de droites représentatives}
Dans un repère orthonormé d'unité $1\text{ cm}$ pour 1 unité sur chaque axe :
\begin{enumerate}
    \item Construire le tableau de valeurs pour la fonction $f(x) = 2x - 3$ avec $x \in \{-1 ; 0 ; 2 ; 4\}$.
    \item Tracer la droite $(d_1)$ représentative de $f$.
    \item Tracer sur le même repère la droite $(d_2)$ représentative de la fonction linéaire $g(x) = -1{,}5x$.
\end{enumerate}
\quadrillage[4.5cm]
\end{exercice}

\begin{exercice}[4 pts]{\capAnalyser}{Problème Commercial : Forfaits de Téléphonie}
Un opérateur propose deux forfaits pour une flotte d'entreprise :
\begin{itemize}
    \item \textbf{Offre A} : Sans abonnement, $0{,}15\euro$ la minute d'appel : $A(x) = 0{,}15x$.
    \item \textbf{Offre B} : Abonnement fixe de $12\euro$ par mois puis $0{,}05\euro$ la minute : $B(x) = 0{,}05x + 12$.
\end{itemize}
\begin{enumerate}
    \item Préciser la nature de chaque fonction (linéaire ou affine).
    \item Calculer le montant de la facture pour 60 minutes de communication avec chaque offre.
    \item Résoudre l'équation $0{,}15x = 0{,}05x + 12$ et interpréter le résultat.
\end{enumerate}
\lignesreponse[5]
\end{exercice}

% ==============================================================================
% 5. VERS LE BREVET (DNB PRO)
% ==============================================================================
\section{Vers le Brevet (DNB Pro)}

\begin{sujetdnb}[DNB Pro Métropole][10 pts]{Gestion d'un Club Nautique}
Un club de kayak propose deux tarifs de location de matériel pour la saison d'été :
\begin{itemize}
    \item \textbf{Tarif Plein} : $18\euro$ par séance de 2 heures.
    \item \textbf{Tarif Club (avec adhésion)} : Carte annuelle à $40\euro$ puis $10\euro$ par séance de 2 heures.
\end{itemize}
On note $x$ le nombre de séances effectuées par un adhérent.

\begin{center}
\begin{tikzpicture}[scale=0.75]
    \draw[step=1cm, slate200, very thin] (0,0) grid (10,10);
    \draw[thick, ->, >=stealth, slate700] (0,0) -- (10.5,0) node[right] {Nombre de séances $x$};
    \draw[thick, ->, >=stealth, slate700] (0,0) -- (0,10.5) node[above] {Prix total (\euro)};

    % Graduations (axe x : 1 car = 1 seance, axe y : 1 car = 20 euros)
    \foreach \x in {1,2,...,10} \draw (\x,0.1) -- (\x,-0.1) node[below, font=\scriptsize] {\x};
    \foreach \y in {20,40,...,200} \draw (0.1,\y/20) -- (-0.1,\y/20) node[left, font=\scriptsize] {\y};

    % Droites
    % Tarif Plein: 18x -> y = 18*x/20
    \draw[line width=1.3pt, colPrimary] (0,0) -- (10,9) node[above right] {$T_1(x) = 18x$};
    % Tarif Club: 10x + 40 -> y = (10*x + 40)/20 = 0.5x + 2
    \draw[line width=1.3pt, colExem] (0,2) -- (10,7) node[above right] {$T_2(x) = 10x + 40$};

    \fill[red] (5,4.5) circle (3pt) node[above left, text=red, font=\footnotesize\bfseries] {Point d'intersection};
\end{tikzpicture}
\end{center}

\begin{enumerate}
    \item \capSApproprier\ Associer à chaque tarif sa fonction représentative $T_1(x)$ et $T_2(x)$.
    \item \capAnalyser\ Par lecture graphique :
    \begin{enumerate}
        \item Quel est le prix à payer pour 3 séances au tarif plein ?
        \item Pour quel nombre de séances les deux tarifs sont-ils rigoureusement identiques ? Quel est alors le prix ?
        \item Pour un jeune qui prévoit de faire 8 séances dans l'été, quel est le tarif le plus avantageux ? Quelle sera son économie ?
    \end{enumerate}
    \item \capRealiser\ Retrouver par le calcul le nombre de séances pour lequel $T_1(x) = T_2(x)$.
\end{enumerate}
\lignesreponse[7]
\end{sujetdnb}

% ==============================================================================
% 6. CORRIGÉS DÉTAILLÉS
% ==============================================================================
\section{Corrigés Détaillés du Chapitre}

\begin{solution}[Corrigé des Exercices et du Sujet DNB]
\textbf{Exercice 1} :
1. $f(2) = 4(2)-7 = 1$ ; $f(-3) = 4(-3)-7 = -19$ ; $f(0) = -7$. \\
2. $4x - 7 = 13 \implies 4x = 20 \implies x = \repbox{5}$. \\
3. $f(5) = 4(5)-7 = 13$. Donc le point $M(5;13)$ \textbf{appartient bien} à la droite.

\vspace{2mm}
\textbf{Exercice 3 (Téléphonie)} :
1. $A$ est une fonction linéaire (proportionnalité), $B$ est une fonction affine. \\
2. $A(60) = 0{,}15 \times 60 = 9\euro$ ; $B(60) = 0{,}05 \times 60 + 12 = 15\euro$. \\
3. $0{,}15x - 0{,}05x = 12 \implies 0{,}10x = 12 \implies x = \repbox{120\text{ minutes}}$. Au-delà de 2 heures d'appel, l'Offre B est plus rentable.

\vspace{2mm}
\textbf{Sujet DNB Pro (Club Nautique)} :
1. $T_1(x) = 18x$ (Tarif Plein) et $T_2(x) = 10x + 40$ (Tarif Club). \\
2. (a) Pour 3 séances au tarif plein : $18 \times 3 = 54\euro$ (graphiquement ordonnée 2,7). \\
(b) L'intersection a lieu pour \textbf{5 séances}, le montant est de \textbf{90 \euro}. \\
(c) Pour 8 séances : $T_1(8) = 144\euro$ et $T_2(8) = 120\euro$. Le **Tarif Club** est plus avantageux, économie de $144 - 120 = \repbox{24\euro}$. \\
3. $18x = 10x + 40 \implies 8x = 40 \implies x = \frac{40}{8} = \repbox{5\text{ séances}}$.
\end{solution}
